\textbf{Цели работы:}
\begin{itemize}
	\item познакомиться с технологией Docker и ее концепциями;
	\item изучить основные команды Docker CLI;
	\item изучить формат описания образов с помощью Dockerfile.
\end{itemize}

\textbf{Задачи:}
\begin{enumerate}
	\item изучить методические материалы к лабораторной работе;
	\item развернуть контейнер nginx:
	\begin{itemize}
		\item с тегом alpine3.20;
		\item с именем <имя пользователя>-my-nginx ;
		\item доступный по порту 8000;
		\item запущенный в фоновом режиме;
	\end{itemize}
	\item поcмотреть логи приложения;
	\item проинспектировать контейнер и c помощью jq вывести:
	\begin{itemize}
		\item время, когда контейнер был создан;
		\item версию nginx;
		\item ip-адрес контейнера;
	\end{itemize}
	\item написать Dockerfile для приложения на Python и сбилдить образ;
	\item написать Dockerfile для приложения на Java:
	\begin{itemize}
		\item без использования multi-stage сборки;
		\item с использованием multi-stage сборки;
		\item сравнить получившиеся образы.
	\end{itemize}
\end{enumerate}

\section*{Выполнение задания}

\begin{figure}[!htbp]
	\centering
	\includegraphics[width=0.6\linewidth]{"Screenshot 2026-04-09 105758"}
	\caption{Установленный Docker}
\end{figure}
\FloatBarrier
%docker desktop start 
%docker --version

\begin{figure}[!htbp]
	\centering
	\includegraphics[width=0.6\linewidth]{"Screenshot 2026-04-09 110126"}
	\caption{Запуск контейнера nginx}
\end{figure}
\FloatBarrier
%docker pull nginx:alpine3.20
%docker run --name yeel-my-nginx -p 8000:8000 -d nginx:alpine3.20

\begin{figure}[!htbp]
	\centering
	\includegraphics[width=0.6\linewidth]{"Screenshot 2026-04-09 110148"}
	\caption{Просмотр логов}
\end{figure}
\FloatBarrier
%docker logs yeel-my-nginx

\begin{figure}[!htbp]
	\centering
	\includegraphics[width=0.6\linewidth]{"Screenshot 2026-04-09 110209"}
	\caption{Инспектирование контейнера}
\end{figure}
\FloatBarrier
%docker inspect yeel-my-nginx | jq '.[].Created'
%docker inspect yeel-my-nginx | jq '.[].NetworkSettings.Networks.bridge.IPAddress'
%docker inspect yeel-my-nginx | jq '.[0].Config.Env[] | select(test(\"NGINX_VERSION\"))'
%docker exec yeel-my-nginx nginx -v

\begin{figure}[!htbp]
	\centering
	\includegraphics[width=0.6\linewidth]{"Screenshot 2026-04-09 110428"}
	\caption{Приветственная страница nginx}
\end{figure}
\FloatBarrier

\begin{figure}[!htbp]
	\centering
	\includegraphics[width=0.6\linewidth]{"Screenshot 2026-04-09 110602"}
	\caption{Сборка образа Python}
\end{figure}
\FloatBarrier
%docker build . -t pyapp
%docker run -p 8080:8080 -d -it pyapp

\begin{figure}[!htbp]
	\centering
	\includegraphics[width=0.6\linewidth]{"Screenshot 2026-04-09 110638"}
	\caption{Сборка образа Java без multi-stage}
\end{figure}
\FloatBarrier
%docker build . -t javaapp:huge -f plain.Dockerfile 

\begin{figure}[!htbp]
	\centering
	\includegraphics[width=0.6\linewidth]{"Screenshot 2026-04-09 110705"}
	\caption{Сборка образа Java c multi-stage}
\end{figure}
\FloatBarrier
%docker build . -t javaapp:slim -f multi-stage.Dockerfile 

\begin{figure}[!htbp]
	\centering
	\includegraphics[width=0.6\linewidth]{"Screenshot 2026-04-09 111251"}
	\caption{Сравнение образов}
\end{figure}
\FloatBarrier
%docker ps -a
%docker images -a

%docker desktop stop
\clearpage

\section*{Исходный код программы}

\lstinputlisting[style=dockerfile]{devops-lab3/python/Dockerfile}
Листинг 1 -- Dockerfile для приложения на Python

\lstinputlisting[style=dockerfile]{devops-lab3/java/plain.Dockerfile}
Листинг 2 -- Dockerfile для приложения на Java без multi-stage сборки

\lstinputlisting[style=dockerfile]{devops-lab3/java/multi-stage.Dockerfile}
Листинг 3 -- Dockerfile для приложения на Java с multi-stage сборкой

\section*{Заключение}

В ходе выполнения лабораторной работы были решены следующие задачи:
\begin{itemize}
	\item проведено ознакомление с технологией Docker и ее концепциями;
	\item изучены основные команды Docker CLI;
	\item изучен формат описания образов с помощью Dockerfile.
\end{itemize}

\section*{Контрольные вопросы}

\begin{enumerate}
	\item \textit{Что такое Docker?} Контейнерная платформа программного обеспечения, созданная для разработки, отправки и запуска приложений с использованием контейнерных технологий..
	\item \textit{Что такое образ?} Схема приложения, которая является основой контейнеров и выступает в роли шаблона, в соответствии с которым создаются контейнеры.
	\item \textit{Что такое контейнер?} Набор изолированных от других процессов, которые создаются на основе образа и запускают само приложение.
	\item \textit{На рисунке ниже продемонстрирован образ, в котором находится один файл main.py. По умолчанию программа выводит в терминал фразу "Hello, World!". Однако есть возможность изменить приветствие, передав в качестве аргумента слово, которое будет напечатано вместо "World". При втором запуске контейнера произошла ошибка. Укажите на причину ошибки и найдите способ ее устранения} 
	\begin{center}
		\includegraphics[width=0.6\linewidth]{"Screenshot 2026-04-09 104859"}
	\end{center}
	Ошибка происходит из-за особенностей обработки аргументов контейнером. Исправить ее можно:
	\begin{enumerate}[label=\alph*.]
		\item заменив CMD ["python", "main.py"] на CMD python main.py (shell-режим);
		\item заменив CMD ["python", "main.py"] на ENTRYPOINT ["python", "main.py"] и CMD [] (exec-режим).
	\end{enumerate}
\end{enumerate}
